\subsection{CM theta normalization and further unconditional deductions}

The holomorphic theta normalization in Lead 16(e) differs from the
norm occurring in the question. We make that distinction explicit,
including the finite-character ambiguity in the multiplication law.

\begin{statement}[C.7.6; Proved: the exact CM multiplication law]
Let $E/\overline{\Q}$ have complex multiplication by an order
$\mathcal O$, and use an algebraic invariant differential. Write
$A=\operatorname{area}(\Lambda)$ and
\[
 \eta(z)=s_2z+\frac{\pi}{A}\bar z,\qquad
 \theta(z)=\sigma(z)e^{-s_2z^2/2}.
\]
Then $s_2\in\overline{\Q}$. For every nonzero
$\alpha\in\mathcal O$, with $d=\alpha\bar\alpha$, there is
$F_\alpha\in\overline{\Q}(E)$ such that
\begin{equation}\label{eq:ct-CMmult}
 \theta(\alpha z)^2=F_\alpha(\exp_E z)\theta(z)^{2d}.
\end{equation}
Its divisor and leading Laurent coefficient are
\[
 \operatorname{div}(F_\alpha)=2[\alpha]^*(O)-2d(O),\qquad
 F_\alpha(\exp_Ez)=\alpha^2z^{2-2d}(1+O(z^2)).
\]
Thus the unsquared quotient is an algebraic function on $E$, but
it need not be a rational function on $E$. For a nontorsion
$P=\exp_Eu\in E(\overline{\Q})$, algebraicity of $\theta(u)$
would imply algebraicity of $\theta(\alpha u)$ for all
$\alpha\in\mathcal O$. The additive group $\mathcal O u$ has
rank two, not more than four.
\end{statement}
\begin{proof}
The Legendre relation gives the displayed real-linear decomposition
of $\eta$. The sigma automorphy formula gives
\[
 \theta(z+\omega)=\epsilon(\omega)
  \exp\left(\frac{\pi\bar\omega}{A}
                         (z+\omega/2)\right)\theta(z),
 \qquad \epsilon(\omega)\in\{1,-1\}.
\]
For $\alpha\Lambda\subset\Lambda$, the exponential multipliers
of $\theta(\alpha z)/\theta(z)^d$ cancel because
$\alpha\bar\alpha=d$. Its remaining multiplier is a sign.
Squaring therefore gives a meromorphic function on $E$, with the
stated divisor. This divisor is defined over $\overline{\Q}$.
The existence of a rational function with this divisor over $\C$
implies its existence over $\overline{\Q}$: the relevant
one-dimensional Riemann--Roch space is obtained by extension of
scalars. Normalize its leading Laurent coefficient, in the formal
parameter of the algebraic differential, to $\alpha^2$. This
identifies it with the squared analytic quotient. The action of
$\alpha$ on this differential is algebraic, so $\alpha^2$ is an
allowed normalization.

Choose a nonreal $\alpha\in\mathcal O$. Since the normalized
sigma series is $\sigma(z)=z+O(z^5)$ with algebraic Laurent
coefficients, the next coefficient of the squared quotient is
\[
 \alpha^2z^{2-2d}
       \{1-s_2(\alpha^2-d)z^2+O(z^4)\}.
\]
The coefficients of the rational function in this algebraic formal
parameter are algebraic. As $\alpha^2-d\ne0$, this proves
$s_2\in\overline{\Q}$ without an unnormalized appeal to a CM
weight-two form. It also proves that the functions in
\eqref{eq:ct-CMmult} are algebraic over the chosen field.
At nontorsion $u$, none of the arguments $\alpha u$ is a lattice
point; evaluating the identity proves the algebraicity implication.
The injective map $\mathcal O\to\C$, $\alpha\mapsto\alpha u$,
proves the rank assertion.

The qualification about rationality is necessary. For a square
lattice, $s_2=0$ and $\alpha=1+i$ give
$F_{1+i}=2i\wp$. The unsquared quotient has a simple zero at
the nonzero point in $\ker(1+i)$ and a simple pole at $O$,
which is not a principal divisor on $E$; its square is the
rational function just displayed.
\end{proof}

\begin{statement}[C.7.7; Known with derivation: holomorphic CM theta values]
Let $E$, $u$ and $\theta$ satisfy the hypotheses of C.7.6, with
$P$ nontorsion. Then $\theta(u)$ and
$\zeta(u)-s_2u$ are transcendental. The first assertion is the
diagonal specialization of Bertrand's Corollary 3, in the exact
form reproduced by its author in
\citet[final slide, comment (c)]{bertrand2017}. It concerns
$\theta$, not $N$.
\end{statement}
\begin{proof}
We state the cited criterion in the needed normalization. For CM
$E/\overline{\Q}$, algebraic nontorsion points $x=\exp_Eu$
and $y=\exp_Ev$, and $u,v,u+v\notin\Lambda$, put
\[
 g(u,v)=\log\frac{\sigma(u+v)}{\sigma(u)\sigma(v)}.
\]
The criterion is that
$g(u,v)-s_2uv\in\log\overline{\Q}^{\times}$ if and only if
$Ny=f(x)$ for some positive integer $N$ and some
Rosati-antisymmetric endomorphism $f$. This is the statement
on the opened primary author's slide; its cited original is
\citet[Corollary 3]{bertrand1984}, whose first two pages, but
not the complete original proof, were accessible in this run.
Our use is of the explicitly reproduced criterion, not an
inference from that paper's title or abstract.

Take $v=u$. The sigma duplication identity yields
\[
 g(u,u)-s_2u^2=2\log\theta(u)+\log(-\wp'(u))
                     \pmod{2\pi i\Z}.
\]
All of the criterion's hypotheses hold: $P$ is algebraic,
$P$ and $2P$ are nonzero, and $\wp'(u)$ is algebraic and
nonzero. If $\theta(u)$ were algebraic, the criterion would
give $(N-f)P=O$. On a CM elliptic curve the Rosati involution
is complex conjugation. Thus $f+\bar f=0$ whereas $N>0$,
so $N-f$ is a nonzero endomorphism. Its kernel is finite,
contradicting nontorsion of $P$.

For the second assertion we apply the analytic subgroup theorem
in \citet[Theorem 6.1]{huber2022}. In the algebraic universal
vector extension $\widetilde E$ a logarithm of an algebraic
lift of $P$ has coordinates $(u,\zeta(u)-b)$, where $b$ is
an algebraic fibre coordinate, with a possible algebraic linear
change of the second coordinate. If
$b=\zeta(u)-s_2u$ were algebraic, this logarithm would belong
to the algebraic line $\{(z,s_2z):z\in\C\}$.
The analytic subgroup theorem applies to the connected
commutative algebraic group $\widetilde E$ over
$\overline{\Q}$ and this line defined over
$\overline{\Q}$; its exponential contains the nonzero
algebraic lift. It would provide a connected one-dimensional
algebraic subgroup projecting nontrivially to $E$. Such a
subgroup is an elliptic curve isogenous to $E$ and would
split the universal vector extension after that isogeny.
This is impossible: its nonzero class in
$\operatorname{Ext}^1(E,\mathbb G_a)=H^1(E,\mathcal O_E)$
remains nonzero under an isogeny, since pullback on this
one-dimensional space is invertible in characteristic zero.
\end{proof}

\paragraph{The archimedean term still requiring control.}
The exact comparison is
\begin{equation}\label{eq:ct-Gaussian}
 N(P)=|\Delta|^{1/12}|\theta(u)|
                         \exp\left(-\frac{\pi|u|^2}{2A}\right).
\end{equation}
It follows from C.7.1 and the real-linear decomposition of
$\eta$. C.7.7 does not exclude algebraicity of this product.
The generic finite-order interpolation approach in Lead 16(e)
supplies only the rank-two set $\mathcal O u$; it does not
meet the proposed rank-greater-than-four criterion. The
derivative approach would need the algebraicity of
$\theta'(u)/\theta(u)=\zeta(u)-s_2u$, which C.7.7 excludes.
Replacing the norm by the holomorphic theta value would
therefore change the problem precisely by the factor in
\eqref{eq:ct-Gaussian}.

We also tested an additive-extension route.
\citet[Theorem 3.1.2]{bost2009} proves that an algebraic point
of the universal vector extension of an abelian variety which
lies in its maximal compact subgroup at one complex embedding
is torsion. The theorem's hypotheses require an algebraic
point of the vector extension, equivalently an algebraic line
bundle with an algebraic connection having unitary monodromy.
C.7.2 supplies only an algebraic point in one fibre of a
multiplicative extension. It does not supply that algebraic
connection. Differentiating a metric value at one algebraic
point cannot establish algebraicity of the connection there
or globally. This is the exact missing implication in that
attempt.

Another route uses the Latt\`es map $f_E$ defined by
$x(2Q)=f_E(x(Q))$. Its canonical height is
$\widehat h_{f_E}(x(P))=2\widehat h_O(P)$. For rational
parameters C.6 therefore makes algebraicity of $N(P)$
equivalent to algebraicity of
$\exp(\widehat h_{f_E}(x(P)))$.
\citet[Theorem 1.4 and Corollary 1.6]{nguyen2023} settle the
corresponding multiplicative-height problem for nonintegrable
polynomials using B\"ottcher coordinates in the basin of
infinity. They do not apply to $f_E$. Indeed an elliptic
duplication Latt\`es map has no exceptional point: the inverse
images of any point contain the images of all successive
division points of an elliptic point, an infinite set. A
polynomial has the exceptional point $\infty$, so no conjugate
of $f_E$ is polynomial. The Julia set is the whole sphere
and the required basin of infinity is absent. This checks
the failed hypothesis, rather than treating a theorem for
polynomials as a theorem for all rational maps.

